\chapter{Perbandingan Senilai}\label{ch:g7:prop}

Kelas 6 bertemu tabel yang \omterm{def:g6:propdata:def}{sebanding} (\cref{ch:g6:propdata}); bab ini
mengubah \omterm{def:g7:prop:table}{perbandingan senilai} menjadi alat yang utuh: aturan perkalian
silang untuk mencari nilai keempat, persen, serta \omterm{def:g7:prop:scale}{skala} peta dan
model. Kelajuan, \omterm{def:g7:prop:table}{perbandingan senilai} yang paling terkenal, dipelajari
pada \cref{ch:g8:speed}.

\section{Mengenali dan melengkapi tabel}

\begin{definition}[Tabel perbandingan senilai]\label{def:g7:prop:table}
Tabel dengan dua baris disebut \emph{tabel perbandingan
senilai}\index{perbandingan senilai} jika bilangan pada baris kedua
adalah bilangan baris pertama yang dikalikan satu bilangan tetap, yaitu
\emph{koefisiennya}. Setara dengan itu: semua \omterm{def:g3:division:remainder}{hasil bagi} kolom
$\frac{\text{bawah}}{\text{atas}}$ bernilai sama.
\end{definition}

\begin{example}\label{ex:g7:prop:recognize}
\begin{center}
\renewcommand{\arraystretch}{1.3}
\begin{tabular}{l|ccc}
atas & $4$ & $10$ & $14$ \\
\hline
bawah & $6$ & $15$ & $21$ \\
\end{tabular}
\end{center}
\omterm{def:g3:division:remainder}{Hasil baginya}: $\frac64 = 1.5$, $\frac{15}{10} = 1.5$,
$\frac{21}{14} = 1.5$: \omterm{def:g6:propdata:def}{sebanding}, dengan koefisien $1.5$.
\end{example}

\begin{theorem}[Aturan perkalian silang]\label{thm:g7:prop:cross}
Pada \omterm{def:g7:prop:table}{tabel perbandingan senilai} dengan kolom $\begin{pmatrix} a \\ b
\end{pmatrix}$ dan $\begin{pmatrix} c \\ d \end{pmatrix}$, kedua \omterm{def:g2:mult:def}{hasil
kali} silangnya sama:
\[
a \times d = b \times c .
\]
Akibatnya nilai keempat dapat dihitung dari ketiga nilai yang lain:
$d = \dfrac{b \times c}{a}$.
\end{theorem}

\begin{proof}
Misalkan $k$ koefisiennya: $b = ka$ dan $d = kc$. Maka
\[
a \times d = a \times kc = k \times ac,
\qquad
b \times c = ka \times c = k \times ac :
\]
kedua \omterm{def:g2:mult:def}{hasil kali} silangnya sama dengan $k \times ac$. Membagi
$ad = bc$ dengan $a$ memberi rumus untuk $d$.
\end{proof}

\begin{example}[Nilai keempat yang sebanding]\label{ex:g7:prop:fourth}
Kalau $7$ buku yang sama beratnya $2.8$ kg, berapa berat $12$ buku
semacam itu? Tabel dan aturan perkalian silang:
\begin{center}
\renewcommand{\arraystretch}{1.3}
\begin{tabular}{l|cc}
buku & $7$ & $12$ \\
\hline
kg & $2.8$ & $m$ \\
\end{tabular}
\end{center}
\[
7 \times m = 2.8 \times 12
\quad\Longrightarrow\quad
m = \frac{2.8 \times 12}{7} = \frac{33.6}{7} = 4.8 \text{ kg}.
\]
Periksa lewat satuan: satu buku beratnya $0.4$ kg, dan dua belas buku $4.8$ kg.
\end{example}

\section{Persen}

\begin{method}[Menerapkan dan mencari persen]\label{met:g7:prop:percent}
\begin{enumerate}
  \item \emph{Terapkan} $t\,\%$: kalikan dengan $\frac{t}{100}$
        ($12\,\%$ dari $350$ adalah $0.12 \times 350 = 42$);
  \item \emph{carilah} persen yang diwakili sebuah bagian: hitunglah
        $\frac{\text{bagian}}{\text{keseluruhan}}$ lalu tulis ulang
        atas $100$ ($21$ dari $60$:
        $\frac{21}{60} = \frac{35}{100} = 35\,\%$);
  \item tanyakan selalu: \emph{persen dari apa?} Keseluruhan yang
        menjadi acuannya itu penting.
\end{enumerate}
\end{method}

\begin{example}\label{ex:g7:prop:percent}
Di sebuah sekolah, $180$ dari $400$ muridnya perempuan. Persennya:
\[
\frac{180}{400} = \frac{180 \div 4}{400 \div 4} = \frac{45}{100}
= 45\,\% .
\]
Sebaliknya: berapa banyak murid laki-laki yang $55\,\%$ itu?
$0.55 \times 400 = 220$. Periksa: $180 + 220 = 400$.
\end{example}

\section{Skala}

\begin{definition}[Skala]\label{def:g7:prop:scale}
Pada peta atau model berskala \emph{skala}\index{skala}
$\frac{1}{n}$ (ditulis $1 : n$), jarak pada peta \omterm{def:g6:propdata:def}{sebanding} dengan
jarak yang sebenarnya, dengan koefisien $\frac1n$:
\[
\text{jarak pada peta} = \frac{1}{n} \times \text{jarak sebenarnya},
\]
keduanya diukur dengan \emph{satuan yang sama}.
\end{definition}

\begin{example}\label{ex:g7:prop:scale}
Pada peta $1 : 50\,000$, dua kota berjarak $7$ cm. Jarak sebenarnya:
\[
7 \times 50\,000 = 350\,000 \text{ cm} = 3\,500 \text{ m} = 3.5
\text{ km}.
\]
Sebaliknya, jalur pendakian sepanjang $12$ km terukur pada peta:
$12$ km $= 1\,200\,000$ cm, jadi $1\,200\,000 \div 50\,000 = 24$ cm.
\end{example}

\begin{omfigure}
\begin{tikzpicture}
\begin{axis}[omaxis,
  width=8.4cm, height=5.6cm,
  xmin=0, xmax=8.6, ymin=0, ymax=34,
  xlabel={peta (cm)}, ylabel={sebenarnya (km)},
  xtick={2,4,6,8}, ytick={10,20,30}]
  \addplot[omDef, thick, domain=0:8.2] {4*x};
  \addplot[omDef, only marks, mark=*, mark size=1.5pt] coordinates
    {(2,8) (4,16) (6,24) (8,32)};
  \draw[dashed, omThm] (axis cs:7,0) -- (axis cs:7,28)
    -- (axis cs:0,28);
\end{axis}
\end{tikzpicture}

{\small \omterm{def:g7:prop:scale}{Skala} adalah \omterm{def:g7:prop:table}{perbandingan senilai}: jarak pada peta melawan
jarak sebenarnya berbaris pada garis lurus yang melalui titik asal (di
sini $1$ cm $\leftrightarrow$ $4$ km; garis putus-putusnya membaca
$7$ cm $\leftrightarrow$ $28$ km).}
\end{omfigure}

\begin{remark}
Grafik adalah uji \omterm{def:g7:prop:table}{perbandingan senilai} yang paling cepat
(\cref{ch:g6:propdata}): titik yang berbaris pada garis lurus
\emph{melalui titik asal} berarti besaran yang \omterm{def:g6:propdata:def}{sebanding}; garis yang
meleset dari titik asal (seperti tarif taksi dengan biaya tetap)
berarti \emph{tidak} \omterm{def:g6:propdata:def}{sebanding}, walaupun ia garis lurus. Fungsi afin,
pada \cref{ch:g9:linfunc}, akan memperjelas hal ini.
\end{remark}

\section{Latihan}

\begin{exercise}[$\star$]\label{exo:g7:prop:1}
Tabel mana yang merupakan \omterm{def:g7:prop:table}{tabel perbandingan senilai}? Sebutkan
koefisiennya kalau ada.
\begin{center}
\renewcommand{\arraystretch}{1.2}
\begin{tabular}{l|ccc}
 & $3$ & $7$ & $11$ \\
\hline
 & $7.5$ & $17.5$ & $27.5$ \\
\end{tabular}
\qquad
\begin{tabular}{l|ccc}
 & $2$ & $5$ & $9$ \\
\hline
 & $6$ & $15$ & $28$ \\
\end{tabular}
\end{center}
\end{exercise}

\begin{exercise}[$\star$]\label{exo:g7:prop:2}
Lengkapilah dengan aturan perkalian silang, sambil menulis
perhitungannya: $5$ kg kentang berharga $6$ euro; berapa harga $8$ kg?
\end{exercise}

\begin{exercise}[$\star$]\label{exo:g7:prop:3}
Sebuah pencetak mencetak $36$ halaman dalam $3$ menit. Dengan laju
yang sama, berapa halaman dalam $5$ menit? Berapa lama untuk $96$ halaman?
\end{exercise}

\begin{exercise}[$\star$]\label{exo:g7:prop:4}
Hitunglah: $30\,\%$ dari $250$; \quad $15\,\%$ dari $60$; \quad
$t\,\%$ sedemikian sehingga $t\,\%$ dari $80$ adalah $20$.
\end{exercise}

\begin{exercise}[$\star$]\label{exo:g7:prop:5}
Dari $25$ lemparan, seorang pemain bola basket memasukkan $18$. Berapa
persen keberhasilannya? (Tulislah ulang \omterm{def:g6:fractions:def}{pecahannya} atas $100$.)
\end{exercise}

\begin{exercise}[$\star$]\label{exo:g7:prop:6}
Pada peta $1 : 25\,000$, sebuah jalan setapak terukur $9$ cm. Berapa
panjangnya yang sebenarnya dalam km? Sebuah danau panjangnya $2$ km: berapa panjangnya pada peta?
\end{exercise}

\begin{exercise}[$\star$]\label{exo:g7:prop:7}
Sebuah mobil model dibuat dengan \omterm{def:g7:prop:scale}{skala} $1 : 43$. Mobil aslinya
panjangnya $4.3$ m. Berapa panjang modelnya, dalam cm?
\end{exercise}

\begin{exercise}[$\star\star$]\label{exo:g7:prop:8}
Sebuah resep memakai $240$ g cokelat untuk $8$ sajian. Alin punya
$300$ g cokelat. Untuk berapa sajian cokelat itu cukup (dalam bilangan
sajian yang bulat)?
\end{exercise}

\begin{exercise}[$\star\star$]\label{exo:g7:prop:9}
Harga sebuah jaket turun dari $80$ menjadi $60$ euro.
\begin{enumerate}
  \item Berapa potongannya dalam euro? Berapa persen dari harga
        semula?
  \item Kemudian harganya naik lagi dari $60$ menjadi $80$ euro.
        Jelaskan mengapa kenaikan itu \emph{bukan} $25\,\%$ ---
        hitunglah persen kenaikan yang benar.
\end{enumerate}
\end{exercise}

\begin{exercise}[$\star\star$]\label{exo:g7:prop:10}
Dua baris sebuah tabel saling \omterm{def:g6:propdata:def}{sebanding}:
\begin{center}
\renewcommand{\arraystretch}{1.2}
\begin{tabular}{l|ccc}
 & $4$ & $x$ & $10$ \\
\hline
 & $y$ & $10.5$ & $17.5$ \\
\end{tabular}
\end{center}
Carilah koefisiennya, lalu $x$ dan $y$.
\end{exercise}

\begin{exercise}[$\star\star\star$]\label{exo:g7:prop:11}
Sebuah mesin fotokopi memperkecil dokumen menjadi $80\,\%$ ukurannya.
Sebuah \omterm{def:g6:lines:objects}{ruas garis} sepanjang $10$ cm difotokopi, lalu \emph{salinannya}
difotokopi lagi dengan setelan yang sama.
\begin{enumerate}
  \item Berapa panjang \omterm{def:g6:lines:objects}{ruas garis} itu setelah satu salinan? Setelah
        dua?
  \item Mengapa jawabannya setelah dua salinan bukan $60\,\%$ ukuran
        aslinya? Persen tunggal yang mana yang bersesuaian dengan dua
        salinan?
\end{enumerate}
\end{exercise}

\section{Soal: Mengukur Bumi dengan sebatang tongkat}

\begin{problem}[{Soal akhir pekan --- bayangan adalah perbandingan
senilai, dan bagaimana Eratosthenes menghitung keliling Bumi pada
tahun 240 SM}]
\label{pb:g7:prop:1}
Dua puluh dua abad yang lalu, tanpa teleskop, tanpa satelit dan
tanpa kalkulator, sang pustakawan Aleksandria mengukur Bumi ---
dengan sebatang tongkat, sebuah sumur, kafilah unta dan matematika
bab ini. Soal ini menelusuri kembali langkahnya: mula-mula
\omterm{def:g7:prop:table}{perbandingan senilai} pada bayangan, lalu perhitungannya yang
termasyhur itu, dan akhirnya bagaimana rupa jawabannya pada \omterm{def:g7:prop:scale}{skala}
manusia.

\medskip
\noindent\textbf{Bagian I --- Bayangan sebatang tongkat.}
Matahari begitu jauh sehingga sinarnya sampai kepada kita saling
\omterm{def:g6:lines:perp}{sejajar}. Pada saat tertentu, semua benda yang tegak dan bayangannya
karena itu saling \omterm{def:g6:propdata:def}{sebanding}.
\begin{enumerate}
  \item Pada tengah hari, sebatang tongkat tegak setinggi $1$ m
        membuat bayangan sepanjang $0.4$ m, dan bayangan sebatang
        pohon terukur $3.2$ m. Berapa tinggi pohon itu
        (\cref{thm:g7:prop:cross})?
  \item Pada saat yang sama, sebuah menara membuat bayangan $26$ m:
        berapa tinggi menaranya? Dan berapa panjang bayangan seorang
        anak setinggi $1.5$ m?
  \item Jelaskan dalam satu kalimat mengapa semua perhitungan ini
        hanya sah pada \emph{saat yang sama} dalam sehari.
  \item Legenda mengatakan Thales mengukur Piramida Besar dengan
        menunggu saat ketika \emph{bayangannya sendiri tepat sepanjang
        tingginya}. Berapa koefisien \omterm{def:g7:prop:table}{perbandingan senilai} pada saat
        itu, dan berapa tinggi piramidanya kalau ujung bayangannya
        berada $147$ m dari titik di tanah tepat di bawah puncaknya?
  \item Rangkumlah Bagian I: pada satu saat tertentu, besaran mana
        yang sama bagi tongkat, pohon, menara dan piramida itu
        (\cref{def:g7:prop:table})?
\end{enumerate}

\medskip
\noindent\textbf{Bagian II --- Perhitungan Eratosthenes.}
Eratosthenes mengetahui dua fakta. Di kota Syene, pada tengah hari
pada titik balik matahari musim panas, matahari berdiri \emph{tepat
di atas kepala}: sinarnya mencapai dasar sumur yang terdalam, dan
tongkat yang tegak tidak membuat bayangan. Di Aleksandria, $800$ km
tepat di sebelah utara, pada saat yang sama, sebatang tongkat tegak
\emph{memang} membuat bayangan, dan sinar mataharinya membentuk sudut
$7.2^\circ$ dengan arah tegak.
\begin{enumerate}[resume]
  \item Jelaskan mengapa kedua pengamatan itu bersama-sama
        membuktikan bahwa tanah di Syene dan tanah di Aleksandria
        tidak saling \omterm{def:g6:lines:perp}{sejajar} --- artinya, permukaan Bumi melengkung.
        (Apa yang akan dilakukan sinar yang \omterm{def:g6:lines:perp}{sejajar} pada dua tongkat
        di Bumi yang datar?)
  \item Pada gambar Bumi yang bulat dengan sinar matahari yang
        \omterm{def:g6:lines:perp}{sejajar}, sudut $7.2^\circ$ di Aleksandria muncul kembali di
        pusat Bumi, sebagai sudut antara arah kedua kota itu. Buatlah
        gambarnya lalu yakinkan dirimu (pembenaran yang rapi memakai
        pasangan sudut yang sama besar pada
        \cref{ch:g7:triangles}).
  \item Berapa bagian dari satu putaran penuh nilai $7.2^\circ$ itu?
  \item Busur dari Syene ke Aleksandria ($800$ km) bersesuaian dengan
        $7.2^\circ$; seluruh \omterm{def:g6:measure:perimeter}{kelilingnya} bersesuaian dengan
        $360^\circ$. Susunlah \omterm{def:g7:prop:table}{perbandingan senilainya} lalu hitunglah
        \omterm{def:g6:measure:perimeter}{keliling} Bumi.
  \item Simpulkan \omterm{def:g4:geometry:circle}{diameter} Bumi, dengan memakai $\pi \approx
        3.14$ (\cref{prop:g6:measure:circle}) dan \omterm{def:g6:decimals:rounding}{pembulatan} ke
        ratusan kilometer terdekat. (Nilai modernnya
        $12\,742$ km --- sedekat apa hasil orang yang bermodal
        tongkat pada tahun 240 SM?)
\end{enumerate}

\medskip
\noindent\textbf{Bagian III --- Bumi pada \omterm{def:g7:prop:scale}{skala} manusia.}
\begin{enumerate}[resume]
  \item Sebuah bola dunia dibuat dengan \omterm{def:g7:prop:scale}{skala} $1 : 40\,000\,000$.
        Dengan \omterm{def:g6:measure:perimeter}{keliling} yang ditemukan pada pertanyaan 9,
        tunjukkanlah bahwa \omterm{def:g6:measure:perimeter}{keliling} bola dunia itu tepat $1$ m.
        Berapa \omterm{def:g4:geometry:circle}{diameternya}, sampai sentimeter terdekat?
  \item Mont Blanc menjulang kira-kira $4.8$ km. Ubahlah $4.8$ km
        menjadi sentimeter, bagilah dengan $40\,000\,000$, lalu
        nyatakan tinggi gunung itu pada bola dunianya dalam
        milimeter. Apa yang dikatakan jawabannya tentang seberapa
        ``benjol'' Bumi yang sesungguhnya?
  \item Seorang pejalan menempuh $40$ km per hari. Dengan laju itu,
        berapa hari untuk seluruh \omterm{def:g6:measure:perimeter}{keliling} Eratosthenes --- dan
        kira-kira berapa tahun itu?
  \item Udara yang dapat dihirup terpusat kira-kira pada $10$ km yang
        pertama di atas tanah. Berapa persen dari \omterm{def:g3:shapes:shapes}{jari-jari} Bumi
        (kira-kira $6\,370$ km) itu
        (\cref{met:g7:prop:percent})? Bulatkan ke persepuluh persen
        terdekat.
  \item Para sejarawan memperkirakan bahwa nilai yang diumumkan
        Eratosthenes, setelah diubah ke satuan modern, mungkin
        kira-kira $42\,000$ km. Dengan mengambil $40\,000$ km sebagai
        nilai yang benar, hitunglah persen kekeliruannya. Simpulkan
        dalam satu kalimat tentang tongkat, sumur dan \omterm{def:g7:prop:table}{perbandingan senilai}.

\end{enumerate}
\end{problem}
